\subsubsection{An all-order zero-lattice formula and a larger theta comparison interval}
\label{subsec:zero-lattice-all-orders}
The spatially dependent heat comparison must be compatible
with two very different phenomena: the increasing density
of known real zeros at positive heat time, and the
possibility of multiple zeros at a transition.
We give exact results in both directions.
We also extend the previously certified all-order
noncentral comparison interval without further numerical
sampling.

\begin{theorem}[All-order Laguerre coefficients at a simple real zero]
\label{thm:zeroelementary}
Let $G$ be a real entire function of genus at most one,
with only real zeros and a simple zero $a$.
List the other zeros, with multiplicities, as
$\{b_j\}_{j\ge1}$ and set
\[
\lambda_j=\frac1{(a-b_j)^2},\qquad
S(a)=\sum_{j\ge1}\lambda_j<\infty.
\]
Let $e_k(\lambda)$ be the elementary symmetric
sum of degree $k$, with $e_0=1$.
Then, for every integer $n\ge1$,
\begin{equation}
\boxed{\quad L_n[G](a)=G'(a)^2 e_{n-1}(\lambda).\quad}
\label{eq:allorderzero}
\end{equation}
If there are infinitely many distinct other zeros, every
$L_n[G](a)>0$ and the exact adjacent-order ratio is
\begin{equation}
\frac{(n+1)(2n+1)L_{n+1}[G](a)}{L_n[G](a)}
=(n+1)(2n+1)\frac{e_n(\lambda)}{e_{n-1}(\lambda)}.
\label{eq:allorderzeroratio}
\end{equation}
Moreover
\begin{equation}
0<\frac{e_n}{e_{n-1}}\le\frac{S(a)}{n},
\qquad
\frac{(n+1)(2n+1)L_{n+1}[G](a)}
     {L_n[G](a)}
\le\frac{(n+1)(2n+1)}n S(a).
\label{eq:allordersymupper}
\end{equation}
\end{theorem}
\begin{proof}
Factor the simple zero out of the canonical product
$G(z)=(z-a)e^{Az+B}\prod_j
(1-z/b_j)e^{z/b_j}$, with the standard
minor modification when a zero is at the origin.
For $z=a\pm iy$ each pair of canonical factors
has product $1+y^2/(a-b_j)^2$ after
normalization at $a$; the linear exponential
terms cancel. Consequently
\[
G(a+iy)G(a-iy)
=G'(a)^2y^2\prod_j(1+\lambda_j y^2).
\]
The product converges locally uniformly because
$\sum_j\lambda_j<\infty$. Equating coefficients
of $y^{2n}$ proves \eqref{eq:allorderzero}.
For a finite family $(\lambda_1,\ldots,\lambda_N)$,
counting how many times each $n$-element subset
appears after first choosing one of its $n$ elements gives
\[
n e_n=\sum_{j=1}^N
\lambda_j\,e_{n-1}(\lambda_{\ne j})
\le \biggl(\sum_j\lambda_j\biggr)e_{n-1}.
\]
Monotone passage to the infinite family proves
\eqref{eq:allordersymupper}. Strict positivity follows
if there are at least $n$ nonzero $\lambda_j$.
\end{proof}

Applied to $G=H_{1/2}$, this theorem describes the
initial-time comparison coefficient at every known simple
real zero. It is a *necessary* pointwise condition,
not a sufficient global height bound: $S(a)$ itself
must be controlled without presupposing zero separation,
and no estimate at nonzero values of $G$ follows
automatically.

\begin{theorem}[Higher-multiplicity obstruction without division]
\label{thm:multizero}
Let $G$ be real entire and suppose $a$ is a real zero
of multiplicity $m\ge2$. Then
\[
L_0[G](a)=\cdots=L_{m-1}[G](a)=0,
\qquad
\boxed{L_m[G](a)=
\left(\frac{G^{(m)}(a)}{m!}\right)^2>0.}
\]
In particular, with $n=m-1$, no finite real
coefficient $c$ can satisfy
\[
(n+1)(2n+1)L_{n+1}[G](a)
\le c\,L_n[G](a).
\]
This is a division-free obstruction at any multiple
real zero, not just the double-zero case.
\end{theorem}
\begin{proof}
Taylor's formula gives
$G(a+iy)=\frac{G^{(m)}(a)}{m!}(iy)^m+O(y^{m+1})$
and the conjugate expression with $-iy$.
Their product is
$\left(G^{(m)}(a)/m!\right)^2y^{2m}
+O(y^{2m+2})$ because it is even in $y$.
Comparing coefficients gives the first assertion.
For $n=m-1$ the left side of the proposed
comparison is strictly positive, whereas its
right side is zero for every finite $c$.
\end{proof}
The endpoint-singular open-slab criterion of
Theorem~\ref{thm:openslab} is compatible with this
obstruction: it does not demand a finite coefficient
\emph{at} the endpoint of the backward heat interval.
A global proof still has to exclude interior collisions,
or replace the pointwise comparison by a controlled
Duhamel defect.

\begin{proposition}[Sharp interval furnished by the proved quadratic moment gate]
\label{prop:largerband}
Let $U_n(s,x)=L_n[H_{1/2-s}](x)$ and
$A_n(s)=U_n(s,0)>0$, $0\le s\le1/2$.
Assume the previously established theta-specific
moment-ratio estimate
\[
0\le r_n(s):=\frac{B_n(s)}{2A_n(s)}<\frac1{48}
\qquad(n\ge1).
\]
Then for every $n\ge1$, every heat time in the slab
and every $|x|<\sqrt{48}$,
\begin{equation}
\boxed{A_n(s)\bigl(1-x^2/48\bigr)
       \le U_n(s,x)\le A_n(s).}
\label{eq:largerbandlower}
\end{equation}
The following explicitly spatial coefficient is
nonnegative, finite on each smaller closed interval,
and satisfies the division-free comparison:
\begin{equation}
\boxed{\quad
c_n(s,x)=
\frac{(n+1)(2n+1)A_{n+1}(s)}
     {A_n(s)(1-x^2/48)},\qquad
k_nU_{n+1}(s,x)\le c_n(s,x)U_n(s,x).\quad}
\label{eq:largerbandcoefficient}
\end{equation}
In particular, throughout the closed strip
$|x|\le13/2$ one may take
\[
c_n(s)=\frac{192}{23}\,
\frac{(n+1)(2n+1)A_{n+1}(s)}{A_n(s)}.
\]
\end{proposition}
\begin{proof}
The correlation representation gives
$|U_n(s,x)|\le A_n(s)$ and,
using $\cos v\ge1-v^2/2$,
\[
U_n(s,x)\ge A_n(s)-\frac{x^2B_n(s)}2
=A_n(s)(1-r_n(s)x^2)
\ge A_n(s)(1-x^2/48).
\]
This lower bound is positive for
$|x|<\sqrt{48}$. Similarly
$k_nU_{n+1}(s,x)\le k_nA_{n+1}(s)$.
Multiplying the former lower bound by the
coefficient in \eqref{eq:largerbandcoefficient}
proves the division-free inequality.
For $|x|\le13/2$,
$1-x^2/48\ge 1-169/192=23/192$,
giving the uniform rational multiplier.
\end{proof}

This extension is a nontrivial improvement on the
previously displayed $|x|\le6$ statement, but not a
global-in-height argument. The apparent singularity
at $|x|=\sqrt{48}$ is a limitation of the quadratic
cosine minorant, not evidence that the genuine
Laguerre coefficient vanishes there.
No analytic-continuation principle transports the
strict positive inequality from this interval to
the whole real line.

\paragraph{Verification scope.}
The finite symmetric-polynomial identities and
the arithmetic band constants are tested in
\path{certificates/xi_zero_lattice_all_orders_exact.py}.
The Lean module
\path{XiKernel/ZeroLatticeAndBand.lean} checks
the abstract multiplicity contradiction and
the exact real-algebraic source comparison.
It does not formalize the infinite canonical product,
the zero density of $H_{1/2}$ or the genuine theta
moment bounds end to end. The missing RH theorem
is the global heat-time/height/all-order source
inequality outside the established central band,
with zero collisions handled independently.
